% ============================================================
% 第四章 Helmholtz 方程的 FFT 求解
% ============================================================
\section{Helmholtz 方程的 FFT 求解}
\label{sec:helmholtz}

\subsection{离散 Helmholtz 方程}

Helmholtz 方程 $(-\nabla^2 + k^2)u = f$ 的五点差分离散为：
\begin{equation}
  \frac{1}{h^2}(4u_{i,j} - u_{i-1,j} - u_{i+1,j} - u_{i,j-1} - u_{i,j+1}) + k^2 u_{i,j} = f_{i,j}
\end{equation}
与 Poisson 方程的唯一区别在于对角元增加了 $k^2$ 项。

在 Fourier 域中，这个修正等价于将特征值从 $\lambda_{p,q}$ 变为 $\lambda_{p,q} + k^2$，
频域求解公式变为：
\begin{equation}
  \hat{u}_{p,q} = \frac{\hat{f}_{p,q}}{\lambda_{p,q} + k^2}
  \label{eq:helm-freq-solve}
\end{equation}

因此，所有 FFT 直接法（FA、CR、FACR、FFT9）均可自然推广到 Helmholtz 方程，
只需在频域除数中加上 $k^2$ 项。

\subsection{共振检测}

当 $k^2$ 等于离散 Laplacian 的某个特征值时，
频域除数 $\lambda_{p,q} + k^2$ 可能为零（或非常小），
导致解不存在或不稳定。

\begin{definition}[共振条件]
  若存在模式 $(p,q)$ 使得 $|\lambda_{p,q} + k^2| < \varepsilon$
  （$\varepsilon$ 为小常数），则称 $k^2$ 接近共振波数。
\end{definition}

本文实现了共振检测函数，在求解前检查频域除数的最小值，
并在检测到共振时发出警告。

\subsection{Neumann 边界条件的 Ghost-Point 对称化方法}

对于 Neumann 边界条件，ghost-point Laplacian 矩阵 $G$ 是非对称的。
如第\ref{sec:math}章所述，通过对角缩放 $S = D^{-1}GD$ 实现对称化，
使得 $S$ 可被 DCT-I 对角化。

\subsubsection{Helmholtz 修正}

对于 Helmholtz 方程，Neumann 边界条件下的频域求解公式为：
\begin{equation}
  \hat{u}_{k,l} = \frac{\hat{f}_{k,l}}{\lambda_k^{\mathrm{Neu}} + \lambda_l^{\mathrm{Neu}} + k^2}
\end{equation}
其中 Neumann 特征值为：
\begin{equation}
  \lambda_k^{\mathrm{Neu}} = \frac{2}{h^2}\left(1 - \cos\frac{k\pi}{n-1}\right),
  \quad k = 0, 1, \ldots, n-1
\end{equation}

注意 $k=0$ 对应 $\lambda_0^{\mathrm{Neu}} = 0$（常数模式），
对于 Poisson 方程（$k^2 = 0$）会导致奇异性，
但对于 Helmholtz 方程（$k^2 > 0$），分母为 $k^2 > 0$，无奇异性。

\subsubsection{CR/FACR 中的 Neumann 处理}

在 CR 和 FACR 方法中，1D 变换后的 $y$-方向三对角系统的结构
依赖于 $x$ 方向的边界条件类型。

对于 $x$-方向 Neumann 边界条件，经过 DCT-I 变换后，
$y$-方向三对角系统的次对角线不再是常数 $-1$，
而是变系数 $\mathbf{e} = [-\sqrt{2}, -1, -1, \ldots, -1, -\sqrt{2}]$。
这要求 CR 和 FACR 使用通用对称三对角 Thomas 算法（thomas\_sym）
而非常数系数版本。

\subsection{混合边界条件}

对于混合边界条件，例如 $x$-方向 Neumann、$y$-方向 Dirichlet，
$x$-方向使用 DCT-I 变换（$n$ 个节点），
$y$-方向使用 DST-I 变换（$n-2$ 个内节点）。
右端项的缩放和边界修正需针对各方向分别处理。
